\chapter{Данные}\label{ch:3}

В работе используются три открытых датасета, каждый из которых исторически специализировался под одну из решаемых задач: DEAM для эмоций, Harmonix Set для структурной разметки и GTZAN для жанровой классификации. Объединение трёх таких разнородных корпусов в общем мультизадачном пайплайне требует унифицированной предобработки и аккуратных сплитов. Оба вопроса разбираются в этой главе.

\section{Используемые датасеты}\label{sec:3-1}

Сводная информация о датасетах приведена в табл.~\ref{tab:datasets}.

\begin{table}[H]
\centering
\caption{Используемые датасеты}
\label{tab:datasets}
\footnotesize
\setlength{\tabcolsep}{4pt}
\begin{tabular}{L{2.3cm}L{2.5cm}L{1.4cm}L{3.2cm}L{1.9cm}L{2.7cm}}
\toprule
Датасет & Размер & Аудио & Аннотации & Задача & Источник \\
\midrule
DEAM & 1802 фрагмента (45~с каждый) & mp3 & Покадровые непрерывные arousal/valence & Эмоции & MediaEval 2013--2018~\cite{aljanaki2017deam} \\
Harmonix Set & 912 треков (полная длительность) & mp3 (загружено 743) & Структурная разметка с временными метками & Сегментация & Nieto и соавт., 2019~\cite{nieto2019harmonix} \\
GTZAN & 999 фрагментов (30~с каждый) & wav & Жанровая метка целиком на трек & Жанр & Tzanetakis \& Cook, 2002~\cite{tzanetakis2002gtzan} \\
Итого & 3713 уникальных файлов & \multicolumn{2}{l}{$\approx 110$~ГБ извлечённых признаков} & 4 задачи & \\
\bottomrule
\end{tabular}
\end{table}

DEAM~\cite{aljanaki2017deam} это один из крупнейших открытых датасетов для распознавания музыкальных эмоций, собранный в рамках конкурсов MediaEval с 2013 по 2018 годы. Корпус включает 1802 фрагмента по 45 секунд из коллекции Free Music Archive (лицензия Creative Commons). Для каждого фрагмента собрана непрерывная покадровая разметка возбуждения и валентности по краудсорсингу: около десяти аннотаторов на трек с шагом 500 миллисекунд. Поставляется он в виде нескольких CSV: усреднённые ряды по 60 фреймам, стандартные отклонения как оценка согласия аннотаторов и общие метаданные. В данной работе непрерывные значения дискретизируются в три категории по порогам $\pm 0.2$: значения ниже $-0.2$ относятся к классу Low (или Dark для валентности), от $-0.2$ до $+0.2$ к классу Mid (Neutral), выше $+0.2$ к High (Bright). Пороги выбирались так, чтобы получить приблизительно равномерное распределение классов (25--40\,\% на класс) и при этом сохранить семантическую осмысленность границ; разбиение по квантилям такой осмысленности не давало.

Harmonix Set~\cite{nieto2019harmonix} это единственный крупный открытый корпус со структурной разметкой коммерческой популярной музыки. Всего в нём 912 треков, для которых вручную в Sonic Visualiser размечены позиции долей, сильных долей и границы секций. На этапе планирования рассматривались также SALAMI и Beatles, но первый смещён в сторону классики и джаза (что плохо согласовалось с двумя другими задачами работы, ориентированными на поп-сцену), а второй слишком мал по объёму. Из всей разметки Harmonix используются только секции. Авторская таксономия включает 17 категорий, что для шестиклассовой постановки избыточно: они сведены к шести базовым меткам. Все вступления (intro, intro\_a, intro\_b) сливаются в Intro; куплеты и темы (verse, verse\_a, verse\_b, theme) в Verse; переходы и переходные секции (bridge, transition, transition\_a) в Bridge; припевы и хуки (chorus, chorus\_a, hook) в Chorus; инструменталы, соло и брейкдауны в Instrumental; концовки и затухания в Outro. Сами аудиозаписи Harmonix Set по понятным причинам не распространяет. Для их скачивания в рамках работы реализован отдельный скрипт (\code{scripts/download\_harmonix\_audio.py}), который по полю \code{youtube\_id} из аннотации загружает аудио через \code{yt-dlp}. Из 912 треков успешно загрузилось 743 (около 81.5\%). Остальные за прошедшие годы удалены с YouTube правообладателями.

GTZAN~\cite{tzanetakis2002gtzan} в работе используется в стандартном виде: 999 рабочих файлов (один оригинальный wav оказался повреждён) разбиты на десять жанров по сто примеров в каждом, длительность одного фрагмента 30 секунд. Жанры стандартные для бенчмарка: blues, classical, country, disco, hip-hop, jazz, metal, pop, reggae, rock. Известные проблемы GTZAN не игнорируются~\cite{sturm2014state}: около 5\,\% файлов представляют собой дубликаты или near-duplicates, у части примеров спорная разметка, есть выраженный artist bias и общая устаревшая таксономия. Несмотря на это, GTZAN сохранён как основной датасет жанровой задачи, поскольку он остаётся единственным практическим способом сопоставимости с двадцатилетней литературой по теме. Поведение модели на современной музыке дополнительно смягчается стратегией top-3 в прототипе (§6.3).

\section{Извлечение признаков}\label{sec:3-2}

Все три датасета унифицируются на этапе извлечения признаков. Из mp3 или wav librosa загружает моноаудио с частотой дискретизации 22050 Гц, после чего считается логарифмическая мел-спектрограмма со стандартными для задач классификации музыки параметрами: 128 мел-фильтров в диапазоне от 30 Гц до частоты Найквиста (11025 Гц), окно STFT в 2048 отсчётов (около 93 миллисекунд) и шаг между фреймами 512 отсчётов (около 23 миллисекунд). Полученные значения мощности логарифмируются через \code{librosa.power\_to\_db(ref=np.max)}, что переводит спектрограмму в шкалу примерно от $-80$ до $0$ децибел.

Полученные спектрограммы сохраняются в файлах формата npz и далее уже не пересчитываются: при каждой эпохе обучения готовые массивы просто читаются с диска. Для скорости это критично. Без такого кэша из-за стоимости декодирования mp3 каждая эпоха стоила бы кратно дороже.

Параметры подытожены в табл.~\ref{tab:logmel-params}.

\begin{table}[H]
\centering
\caption{Параметры логарифмической мел-спектрограммы}
\label{tab:logmel-params}
\begin{tabular}{L{3cm}L{3cm}L{8cm}}
\toprule
Параметр & Значение & Обоснование \\
\midrule
Sample rate & 22050~Гц & Стандарт librosa, покрывает до 11025~Гц \\
\code{n\_fft} & 2048 & STFT окно около 93~мс \\
\code{hop\_length} & 512 & Шаг между фреймами около 23~мс \\
\code{n\_mels} & 128 & Стандарт для классификации музыки \\
\code{fmin} & 30~Гц & Нижняя граница мел-фильтров \\
\code{fmax} & 11025~Гц & Частота Найквиста \\
\bottomrule
\end{tabular}
\end{table}

Затем спектрограмма каждого трека режется на окна одинаковой длительности. В работе используются окна по 10 секунд (431 спектральный фрейм при выбранных параметрах) с шагом 1 секунда между соседними окнами. Длительность 10 секунд выбрана по двум причинам. Первая: этого хватает на типичный музыкальный период (куплет или припев целиком), что критично для структурной задачи. Вторая: AST предобучен ровно на десятисекундных окнах AudioSet, и сохранение этой длительности позволяет использовать предобученные позиционные эмбеддинги без интерполяции. Шаг в одну секунду даёт около пятидесяти окон на минуту записи, чего более чем достаточно для последующего декодирования Витерби.

Сразу после нарезки выполняется глобальная нормализация. По обучающему сплиту вычисляются средние значения и стандартные отклонения отдельно по каждой мел-полосе, сохраняются в \code{stats.npz}, а далее применяются ко всем сплитам: на каждом окне из каждого мел-фильтра вычитается соответствующее среднее и делится на стандартное отклонение. Тот же \code{stats.npz} используется и при инференсе пользовательских аудиозаписей: \code{musicml/infer.py} загружает его автоматически. Без этого шага на инференсе наблюдался характерный сбой: блюзовая запись с высокой долей низких частот ошибочно распознавалась как hip-hop из-за того, что её распределение энергии по фильтрам оказывалось сдвинуто относительно нормализованного обучающего распределения.

\section{Стратегия разбиения на обучающую, валидационную и тестовую выборки}\label{sec:3-3}

Каждый из трёх датасетов разбивается независимо в соотношении 70/15/15 с фиксированным сидом 42. Главное принципиальное решение здесь, ради которого пришлось аккуратно реализовать сплитер, состоит в том, что разбиение делается на уровне треков, а не на уровне окон. Если резать окна и распределять их случайно, окна одного и того же трека неизбежно попадают и в обучающую, и в тестовую выборки, что искусственно завышает метрики обобщения.

Для GTZAN стратификация по жанрам делается автоматически (по сто треков на каждый из десяти классов). Для DEAM стратификация по эмоциональным классам невозможна в принципе: метки покадровые и непрерывные, а разные фрагменты одного трека вполне могут оказаться в разных классах после квантизации, поэтому используется случайное разбиение. Harmonix тоже разбивается случайно: жанровая информация в нём отсутствует, а стратифицировать по структуре бессмысленно, поскольку все треки покрывают все секции. Списки треков каждого сплита сохраняются в \code{splits.json} для воспроизводимости любого запуска.

\section{Дисбаланс классов}\label{sec:3-4}

Из четырёх рассматриваемых задач три (возбуждение, валентность, жанр) сбалансированы практически по построению или близки к этому, и для них специальной обработки не требуется. Иначе обстоит дело со структурной сегментацией, где классы распределены сильно неравномерно: куплет и припев в сумме занимают около 60\% общего времени обучающих треков, тогда как на концовку приходится меньше 10\%. Распределение классов сведено в табл.~\ref{tab:class-distribution}.

\begin{table}[H]
\centering
\caption{Распределение классов в обучающих сплитах}
\label{tab:class-distribution}
\begin{tabular}{L{3cm}L{3cm}L{3cm}L{4cm}}
\toprule
Задача & Класс & Доля окон & Вес inverse-frequency \\
\midrule
Segment & Verse & $\sim$35\% & 0.71 \\
 & Chorus & $\sim$25\% & 1.00 \\
 & Bridge & $\sim$12\% & 2.08 \\
 & Intro & $\sim$10\% & 2.50 (обрезан) \\
 & Instrumental & $\sim$10\% & 2.50 (обрезан) \\
 & Outro & $\sim$8\% & 2.50 (обрезан) \\
Arousal & Low / Mid / High & $\sim$33\% / 34\% / 33\% & $\approx 1.00$ каждый \\
Valence & Dark / Neutral / Bright & $\sim$33\% / 33\% / 34\% & $\approx 1.00$ каждый \\
Genre & любой из 10 жанров & $\sim$10\% каждый & $\approx 1.00$ каждый \\
\bottomrule
\end{tabular}
\end{table}

Без специальной обработки дисбаланса модель склонна к вырожденному решению с преимущественным предсказанием доминирующих классов. Чтобы этого избежать, в работе применяются сразу три механизма, действующие на разных уровнях. На уровне функции потерь применяется focal loss~\cite{lin2017focal} с показателем $\gamma = 1.5$: множитель $(1 - p_t)^\gamma$ уменьшает вклад тех примеров, по которым модель уже уверена и обычно права, перенаправляя градиент на трудные случаи. На уровне отдельных классов используются веса, обратно пропорциональные частотам, но с верхним порогом 2.5. Без такого порога редкие классы получали бы слишком большой вес и провоцировали обратную проблему. Наконец, на уровне общей мультизадачной потери головам сегментации присваивается дополнительный вес 1.5, тогда как остальным задачам по 1.0.

Эмпирически эти три ограничения находятся в довольно узком оптимуме. В одном из экспериментов (CNN v3, см. §5.5) была проверена гипотеза об одновременном усилении всех трёх механизмов: вес головы сегментации до 3.5, $\gamma$ до 2.5, верхний порог весов классов до 3.5. Модель в результате переобучилась на редкие классы и потеряла около 13 процентных пунктов на сегментной задаче. Эта неудача стала одним из главных аргументов в пользу того, что в задаче компенсации дисбаланса дальнейшее усиление веса редких классов приводит к обратному переобучению.
